\documentclass[11pt]{article}

\usepackage[margin=1in]{geometry}
\usepackage{graphicx}
\usepackage{amsmath,amssymb}
\usepackage{booktabs}
\usepackage{hyperref}
\usepackage{natbib}
\usepackage{siunitx}

\title{Is a structural-MRI foundation model worth its weight?\\
  Ridge regression on T1Prep morphometry as a baseline on the
  Dallas Lifespan Brain Study}
\author{Morgan Hough \\
  \texttt{morgan.hough@gmail.com}}
\date{Draft --- \today}

\begin{document}
\maketitle

\begin{abstract}
The release of the MedARC \texttt{smri-fm} dataset, the publication of
BrainIAC \citep{brainiac2026}, and the MICCAI 2025 FOMO25 and SSL3D
challenges won by \citet{gordaliza2026} have renewed interest in
structural-MRI foundation models (sMRI-FMs) as a general-purpose
backbone for neuroimaging. Whether FMs provide a meaningful advance
over classical morphometric features for core tasks --- above all,
brain-age prediction --- has not been settled. In fact
\citet{schulz2025} show that for \emph{disease detection}, simpler
models beat complex ones despite losing on age MAE: ridge regression
with 1\,400 parameters produces larger patient-vs-control effect sizes
than a 46M-parameter CNN on the same data. We therefore construct
a \textbf{four-rung classical morphometry ladder} on the Dallas
Lifespan Brain Study (DLBS, \citealp{park2025dlbs}) and benchmark
three sMRI foundation-model checkpoints against it --- BrainIAC
\citep{brainiac2026}, the FOMO25 challenge-winning
MultiModalUNetVAE \citep{gordaliza2026,fomo25release}, and
AnatCL \citep{barbano2024anatcl} --- under identical bias-correction
regimes. DLBS is, to our knowledge, the largest fully-longitudinal,
fully-public adult lifespan T1 set, making it the natural venue to
test the Vidal-Pi\~neiro critique \citep{vidalpineiro2021} that
cross-sectional brain-age gap is uncorrelated with longitudinal brain
change. We pre-register four analyses: (i) cross-sectional brain-age
MAE per rung and per FM; (ii) patient-vs-control effect size per
\citet{schulz2025} on a matched held-out set; (iii) longitudinal
stability of brain-age gap across the three DLBS waves; and
(iv) whether the FM arms meaningfully beat classical morphometry on
any of these dimensions.
\end{abstract}

\section{Introduction}

Brain-predicted age, introduced by \citet{franke2010} and popularised by
Cole and colleagues \citep{cole2017}, has become the single
most-reported biomarker derived from structural MRI. A decade of
architectural progress --- from relevance-vector regression on VBM
features, through Gaussian-process regression
\citep{cole2017,brainager} to 3D CNNs \citep{peng2021sfcn} --- has pushed
UK Biobank mean absolute error (MAE) from $\sim 5$ years to $\sim 2$ years
\citep{peng2021sfcn}. At the same time, multiple authors have shown that
classical ridge- or kernel-based regressors on FreeSurfer or CAT12
morphometric features match deep networks within noise
\citep{more2023,lombardi2021,bashyam2020}. The apparent ceiling is set by
label noise and site effects, not model capacity.

The sMRI-FM generation --- \citet{brainiac2026} and the
MedARC \texttt{smri-fm} dataset
\footnote{\url{https://huggingface.co/datasets/medarc/smri-fm}} --- trades
architectural novelty for a different value proposition: one backbone,
many downstream tasks. The relevant question is therefore not
``is the FM better?'' but ``is FM + ridge better than
T1Prep + ridge?'' under the same bias-corrected evaluation. If not, the
FM's value lies elsewhere (low-label transfer, rare-disease
classification) and brain-age benchmarks stop being a useful selling
point for the foundation-model framing.

\section{Background}

\subsection{Brain age: origins and the classical baseline}

\citet{franke2010} built the first published brain-age model on
voxel-based morphometry (VBM) features produced by the VBM8/CAT toolbox
--- features directly analogous to what T1Prep produces today. Their
relevance vector regressor achieved \SI{5}{years} MAE on healthy adults,
and set the template for a decade of VBM-plus-classical-regressor
work. Cole's \texttt{brainageR} \citep{brainager} is the canonical
open implementation.

\subsection{The deep-learning era (2017--2022)}

\citet{cole2017} demonstrated that 3D CNNs on minimally-preprocessed
T1s match and slightly exceed VBM-plus-classical pipelines on healthy
UK Biobank data, and produce heritable residuals. \citet{peng2021sfcn}
formalised this with SFCN, a VGG-style 3D net trained with a
soft-label KL-divergence loss
\begin{equation}
  \mathcal{L}_{\text{SFCN}}
  = \mathrm{KL}\left(
      \mathcal{N}(y, \sigma^2) \,\Vert\, \hat{p}(\cdot \mid x)
    \right),
\end{equation}
reporting \SI{2.14}{years} MAE on held-out UK Biobank.
\emph{The paper itself notes} that a post-hoc linear bias correction on
the SFCN outputs maintains state-of-the-art, and that with $N \approx 50$
training subjects deep learning barely edges out linear baselines. The
gap has not meaningfully widened since.

\subsection{The ridge counter-literature}

\citet{lombardi2021} reported \SI{2.7}{years} MAE on Dominantly
Inherited Alzheimer Network (DIAN) data using ridge regression on
FreeSurfer features. \citet{more2023} ran one of the largest systematic
comparisons to date and concluded that kernel-ridge and GP regressors on
morphometric summaries generalise better across sites than 3D
CNNs. The emerging consensus --- rarely stated that sharply in print ---
is that \emph{the ceiling on healthy cross-sectional brain age is set by
label and site noise, not by model capacity.}

This has direct implications for the sMRI-FM value proposition:
if the embedding head is itself a ridge regressor (as it typically
is in frozen-backbone evaluations), its advantage over T1Prep-feature
ridge reduces to whether the FM representation captures age-relevant
variance that classical morphometry misses.

\subsection{Regression dilution and bias correction}

Every brain-age model exhibits an age-dependent residual: over-prediction
for young subjects, under-prediction for old ones. The phenomenon is a
mathematical consequence of minimising MSE under a noisy label
distribution and is formally a regression-to-the-mean artefact
\citep{smith2019, liang2019}. Three correction families dominate
current practice:

\paragraph{Cole / Smith linear correction.} Regress $\hat{y}$ on $y$,
recover $\hat{y} = \alpha y + \beta$, and transform
\begin{equation}
  \hat{y}_{\text{corr}} = \frac{\hat{y} - \beta}{\alpha}.
\end{equation}

\paragraph{Beheshti correction.} Regress $\mathrm{PAD} = \hat{y} - y$
directly on $y$ and subtract the fit.

\paragraph{Age-level (Zhang) correction.} Local $z$-score within each
chronological age bin:
\begin{equation}
  \mathrm{PAD}^{\text{ac}}_i
  = \frac{\mathrm{PAD}_i - \mu_a}{\sigma_a}.
\end{equation}
\citet{zhang2023agelevel} show global linear corrections leave
systematic within-age residuals.

\citet{delange2022mindthegap} is the canonical cautionary review: any
downstream association with phenotype $X$ that is itself age-dependent
can arise purely through uncorrected regression dilution, and reporting
uncorrected BAG is no longer defensible.

\subsection{The longitudinal critique}

\citet{vidalpineiro2021} deliver the most uncomfortable finding in the
field: cross-sectional BAG is \emph{essentially uncorrelated} with the
within-subject rate of brain change measured longitudinally. They argue
that high BAG indexes early-life differences (birth weight, polygenic
scores) more than ongoing neurodegeneration. The implication: a meaningful
evaluation of any brain-age model, FM-derived or not, requires
longitudinal data with repeat scans per subject, and should report
longitudinal $\Delta$BAG in addition to cross-sectional MAE. DLBS ---
three waves, up to $\sim 13$ years separation --- is well suited.

\subsection{The foundation-model era (2024--2026)}

BrainIAC \citep{brainiac2026} is the first published sMRI-FM with
credible generalist claims: SimCLR pretraining on $\sim 48{,}500$ scans
across 35 datasets, validated on seven downstream tasks including brain
age, MCI classification, IDH mutation, and survival. Preprocessing is
deliberately minimal: N4 bias correction, $1\,\mathrm{mm}$ isotropic
resampling, rigid MNI registration, HD-BET skull strip, crop to
$128^3$. Preliminary evaluation on 100 DLBS images (this project,
2026-04-21) shows BrainIAC's predictions concentrate in a narrow
younger band --- consistent with distribution-shift between BrainIAC's
training corpus and the DLBS older-adult cohort.

\citet{barbano2024anatcl} (AnatCL) built a weakly-contrastive SSL
model using cortical thickness, grey-matter volume, and surface area
per Desikan--Killiany region plus chronological age as explicit
anatomical supervision, achieving \SI{2.61}{years} MAE on OpenBHB with
only $3{,}984$ training subjects --- the cleanest existing argument
that anatomical-prior-aware pretraining is data-efficient.

\citet{gordaliza2026} (arxiv 2601.13166), released as the
\texttt{MultiModalUNetVAE} (mmunetvae) checkpoint trained on
FOMO-60K (11\,187 subjects, MAE-style SSL with $4\times 4 \times 4$
voxel masking on $96^3$ patches), won MICCAI 2025 in both SSL3D and
FOMO25 brain-MRI foundation-model challenges. Crucially, their
architecture is a 3D U-Net with VAE --- 10$\times$ smaller and 1--2
orders of magnitude faster to train than transformer-based
competitors. This is the strongest empirical argument to date that
3D U-Net MAE is at least as effective as 3D ViT MAE for brain MRI,
and reframes the sMRI-FM architecture-choice debate.

The MedARC \texttt{smri-fm} dataset
\footnote{\url{https://huggingface.co/datasets/medarc/smri-fm}} is a
community re-implementation in this lineage currently under active
development; the group's prior cortex-MAE on fMRI flat maps
\citep{medarcsmrifm2026} strongly suggests MAE (not SimCLR) will be
their pretraining objective of choice.

Three questions the published literature has not yet answered:
\begin{enumerate}
  \item Does FM + ridge beat morphometry + ridge on brain age, under
        identical bias correction?
  \item Does FM-derived BAG pass the Vidal-Pi\~neiro longitudinal
        stability test, or does it inherit the same disconnect between
        cross-sectional and longitudinal signals?
  \item Per \citet{schulz2025}, does the FM's accuracy-vs-effect-size
        trade-off favour it over a simple ridge on classical
        morphometric features for clinically-relevant phenotype
        detection?
\end{enumerate}

\section{Methods (pre-registration)}

\subsection{Data}

We use the Dallas Lifespan Brain Study \citep{park2025dlbs}, OpenNeuro
\texttt{ds004856}, with 464 / 338 / 224 subjects across waves 1/2/3
(total $\sim 1000$ T1 sessions, Philips 3T Achieva, $\sim 4$--$5$ year
intervals). Ages span \numrange{20}{90} years at wave 1. Second
held-out: ds003592 \citep{spreng2022}, Spreng et al.\ neurocognitive
aging, for external generalisation. Both downloaded in full from
OpenNeuro S3 on 2026-04-22.

\subsection{Preprocessing arms and four morphometry rungs}

For each subject we produce four independent classical feature sets
(rungs) and three foundation-model embeddings. All pipelines consume
the same raw DLBS T1 (FLAIR where required). For reproducibility each
pipeline runs inside a pinned Docker image; builds and digests are
recorded in \texttt{comparison/containers/IMAGE\_DIGESTS.md}.

\paragraph{Rung 1 --- SynthSeg volumes (MedARC baseline).}
The MedARC \texttt{preprocessing/pipeline.py} as of commit
\texttt{fbdb1a83}: SynthStrip brain extraction, ANTs rigid
registration to \texttt{MNI152NLin2009cAsym} at \SI{1}{mm} isotropic,
\texttt{mri\_synthseg --parc --robust}. Produces 33 whole-brain
regions + 34$\times$2 Desikan--Killiany cortical regions + the
Bethlehem five summary metrics GMV/WMV/sGMV/VentCSF/TCV
\citep{bethlehem2022}. Following Nima's 2026-04-21 Discord report we
aggregate to a 71-feature ICV-normalised vector.

\paragraph{Rung 2 --- SynthSeg + N4.} Same as rung 1 but with N4 bias
correction via SimpleITK applied before SynthSeg. Tests Connor's
2026-04-13 proposal empirically and settles whether adding a classical
bias-corrector upstream of a bias-robust network helps, hurts, or
ties.

\paragraph{Rung 3 --- FastSurfer VINN.} Replaces SynthSeg with
FastSurfer's \texttt{aparc.DKTatlas+aseg.deep} VINN segmentation
together with CerebNet (fine cerebellar parcellation) and HypVINN
(hypothalamic subregions). Same ICV normalisation. Isolates the
effect of seg-network quality from bias handling.

\paragraph{Rung 4 --- T1Prep thickness + Jacobian.} T1Prep v0.3.3
(DeepMriPrep + AMAP + CAT-Surface). Extracts (i) AMAP-derived
GM/WM/CSF tissue probability maps warped to MNI152 via Geodesic
Shooting, (ii) cortical thickness per Desikan--Killiany ROI from the
CAT-Surface pial/white reconstruction, (iii) log-Jacobian summary of
the nonlinear warp per coarse parcellation, (iv) total intracranial
volume. Adds the two feature axes (thickness, Jacobian) that volumetric
segmentation cannot capture.

\subsection{Models}

For each of the four morphometry rungs we fit a ridge regressor with
leave-one-site-out cross-validation (DLBS has one site, so 5-fold CV
within site and external validation on ds003592). For each of the
three FM arms we extract embeddings on MedARC-pipeline-preprocessed
input and fit a ridge head on the embedding:

\begin{itemize}
  \item \textbf{Ridge / kernel-ridge} per rung.
  \item \textbf{Gaussian Process} with Mat\'ern-3/2 kernel (rungs 1--4)
        as a second classical reference.
  \item \textbf{FM + ridge:} BrainIAC \citep{brainiac2026},
        FOMO25 mmunetvae \citep{gordaliza2026}, AnatCL
        \citep{barbano2024anatcl}.
  \item \textbf{Fine-tuned SFCN} \citep{peng2021sfcn} as the
        supervised deep-learning reference.
\end{itemize}

All models trained on wave-1 DLBS ($N \approx 464$) with 5-fold CV;
wave-2 and wave-3 held out for longitudinal analysis, ds003592 held
out for external-site generalisation.

\subsection{Bias correction}

For every model we report MAE under four regimes: raw, Cole/Smith,
Beheshti, and Zhang age-level. Downstream associations use the
\citet{smith2019} covariate approach ($y + y^2$) in addition to
corrected BAG.

\subsection{Longitudinal evaluation}

For each returning subject we compute $\Delta\mathrm{BAG}_{i,t}$
between waves and test (a) within-subject stability (test-retest ICC
across waves, after age correction) and (b) correlation with change in
T1Prep-derived morphometric summaries (cortical thickness, ventricle
volume). This is the Vidal-Pi\~neiro test applied to FM outputs for the
first time.

\section{Expected outcomes}

We pre-register four specific hypotheses:

\begin{enumerate}
  \item \emph{Cross-sectional ceiling.} FM + ridge and
        rung-4 (T1Prep thickness + Jacobian) + ridge will match
        within \SI{0.3}{years} MAE on wave 1 held-out subjects.
        Preliminary evidence: Nima's rung-1 ridge baseline on 100
        DLBS achieves MAE $= \SI{6.66}{years}$ uncorrected;
        \citet{lombardi2021} report \SI{2.7}{years} with a better
        feature set --- both below BrainIAC's preliminary DLBS
        behaviour, which underpredicts by design on elderly cohorts.
  \item \emph{Bias persistence.} The regression-dilution slope
        $\alpha$ before correction will be statistically
        indistinguishable across preprocessing arms and across
        FM arms, confirming the bias originates in the label
        distribution and is not representation-specific.
  \item \emph{Longitudinal disconnect.} Following
        \citet{vidalpineiro2021}, within-subject $\Delta\mathrm{BAG}$
        will correlate weakly ($|r| < 0.2$) with within-subject
        longitudinal morphometric change, for FM \emph{and} each
        morphometry rung.
  \item \emph{Disease-effect-size inversion.} Following
        \citet{schulz2025}, we expect rungs 1--4 ridge regressors to
        have \emph{higher} patient-vs-control effect size than any
        FM + ridge head on an AD-spectrum labelling of DLBS, despite
        losing on age MAE. This is the paper's core claim: for the
        clinically interesting outcome, the FM is a regression.
\end{enumerate}

A fifth, exploratory question: does T1Prep's nonlinear MNI warp and
cortical-surface output, if fed to the FM as an alternative input,
change the FM's longitudinal or disease-effect-size behaviour?

\section{Discussion (placeholder)}

To be written after results. The paper's contribution is not a new
model but an apples-to-apples evaluation that pins down where the
sMRI-FM proposition is and is not load-bearing.

\bibliographystyle{plainnat}
\bibliography{refs}

\end{document}
